\documentclass[12pt,a4paper]{article}
\usepackage{iftex}
\ifPDFTeX
  \usepackage[utf8]{inputenc}\usepackage[T2A]{fontenc}\usepackage[russian]{babel}
\else
  \usepackage{fontspec}
  \setmainfont{Times New Roman}\setsansfont{Arial}\setmonofont{Courier New}
  \usepackage[russian]{babel}
\fi
\usepackage{amsmath,amssymb}
\usepackage[margin=2.2cm]{geometry}
\numberwithin{equation}{section}
\begin{document}

\section{Поле и потенциал дефлектора}

В плоскости кольца поле записываем до второго порядка:
\begin{equation}\label{eq:Ex}
E_x=E_0\Bigl(1-n\frac{x}{\rho}+m\frac{x^2}{\rho^2}+O(x^3)\Bigr).
\end{equation}
Для поля $E\propto r^{-n}$, $r=\rho+x$, коэффициент $m=n(n+1)/2$: $m=1$ для цилиндрического
($n=1$) и $m=3$ для сферического ($n=2$) дефлектора. В общем случае $m$ --- независимый параметр
формы электродов.

Как и раньше, $qE_0=-p_sv_s/\rho$ и $\chi=q\Phi/(p_sv_s)$. Интегрируя $E_x=-\partial\Phi/\partial x$ при $y=0$,
\begin{equation}\label{eq:chi0}
\chi(x,0)=u-\frac n2u^2+\frac m3u^3,\qquad u=\frac x\rho .
\end{equation}

\paragraph{Уравнение Лапласа.} Между электродами зарядов нет, $\nabla\cdot\vec E=0$, и поле
потенциально, $\vec E=-\nabla\Phi$, поэтому $\nabla^2\Phi=0$. В координатах дефлектора
$(x,y,s)$ с масштабным множителем $h_s=1+x/\rho$
\begin{equation}
\nabla^2\Phi=\frac{1}{h_s}\left[\frac{\partial}{\partial x}\Bigl(h_s\frac{\partial\Phi}{\partial x}\Bigr)
+\frac{\partial}{\partial y}\Bigl(h_s\frac{\partial\Phi}{\partial y}\Bigr)
+\frac{\partial}{\partial s}\Bigl(\frac{1}{h_s}\frac{\partial\Phi}{\partial s}\Bigr)\right].
\end{equation}
Поле от $s$ не зависит, $\partial h_s/\partial x=1/\rho$, и
\begin{equation}\label{eq:laplace}
\frac{\partial^2\chi}{\partial x^2}+\frac{\partial^2\chi}{\partial y^2}
+\frac{1}{\rho+x}\frac{\partial\chi}{\partial x}=0
\end{equation}
--- лапласиан в цилиндрических координатах с $r=\rho+x$.

Поле в плоскости $y=0$ не задаёт зависимость от $y$. Ищем
\begin{equation}
\chi=u-\frac n2u^2+\frac m3u^3+a\,v^2+b\,uv^2,\qquad v=\frac y\rho .
\end{equation}
Подставляя в \eqref{eq:laplace}, получаем в нулевом порядке $-n+2a+1=0$, в первом по $u$:
$2m+2b-(n+1)=0$. Отсюда
\begin{equation}\label{eq:chi}
\boxed{\chi=\frac x\rho-\frac n2\frac{x^2}{\rho^2}+\frac{n-1}{2}\frac{y^2}{\rho^2}
+\frac m3\frac{x^3}{\rho^3}+\Bigl(\frac{n+1}{2}-m\Bigr)\frac{xy^2}{\rho^3}}
\end{equation}
Для цилиндра ($n=m=1$) члены с $y$ пропадают, как и должно быть.


\subsection{Общий вид потенциала}

\paragraph{Ряд и рекуррентное соотношение.} Запишем безразмерный потенциал рядом
\begin{equation}\label{eq:series}
\chi(u,v)=\sum_{i\ge0,\;j\ \text{чётн.}}a_{ij}\,u^iv^j,\qquad u=\frac x\rho,\quad v=\frac y\rho,
\end{equation}
нечётных степеней $v$ нет из-за симметрии относительно медианной плоскости. Умножив
\eqref{eq:laplace} на $1+u$, получаем
$\partial_u\bigl[(1+u)\partial_u\chi\bigr]+(1+u)\partial_v^2\chi=0$, и при $u^iv^j$
\begin{equation}\label{eq:rec}
\boxed{(j+2)(j+1)\bigl(a_{i,j+2}+a_{i-1,j+2}\bigr)
=-(i+2)(i+1)\,a_{i+2,j}-(i+1)^2\,a_{i+1,j}}
\end{equation}
($a_{-1,j}=0$). Все коэффициенты с $v^2,v^4,\dots$ последовательно выражаются через
коэффициенты $a_{i0}$ в медианной плоскости, т.\,е. через $E_x(x,0)$. Поэтому достаточно
задать поле на медианной плоскости:
\begin{equation}\label{eq:Ex-gen}
E_x(x,0)=E_0\Bigl(1-n\,u+\frac{K_2}{2!}u^2+\frac{K_3}{3!}u^3+\dots\Bigr),
\end{equation}
$n$ --- показатель поля, $K_2,K_3,\dots$ --- высшие коэффициенты (в \eqref{eq:Ex} $m=K_2/2$).
Тогда $a_{10}=1$, $a_{20}=-n/2$, $a_{30}=K_2/6$, $a_{40}=K_3/24,\dots$, и из \eqref{eq:rec}
\begin{equation}
a_{02}=\frac{n-1}{2},\qquad a_{12}=\frac{n+1}{2}-\frac{K_2}{2},
\end{equation}
что совпадает с \eqref{eq:chi}.

\paragraph{Элементарные решения.} Решения уравнения Лапласа с осевой симметрией относительно
оси кольца, чётные по $y$, в сферических координатах с центром в центре кольца --- это
$R^{l}P_l(\cos\theta)$, $R^{-(l+1)}P_l(\cos\theta)$ с чётным $l$, а в цилиндрических ---
ещё $\ln r$. Нормированные так, что $\chi=u+\dots$ на орбите, три простейших:
\begin{equation}\label{eq:basis}
\begin{aligned}
\chi_{\rm sph}&=1-\frac{1}{\sqrt{(1+u)^2+v^2}}, & n&=2,\ K_2=6,\\
\chi_{\rm cyl}&=\ln(1+u), & n&=1,\ K_2=2,\\
\chi_{\rm hyp}&=\tfrac12\bigl[(1+u)^2-2v^2-1\bigr], & n&=-1,\ K_2=0
\end{aligned}
\end{equation}
--- сферический, цилиндрический и гиперболоидальный дефлекторы. Каждая функция точно
удовлетворяет \eqref{eq:laplace}, и её разложение имеет вид \eqref{eq:chi}.

\paragraph{Линейные комбинации.} Сумма решений --- тоже решение. Для
$\chi=c_s\chi_{\rm sph}+c_c\chi_{\rm cyl}+c_h\chi_{\rm hyp}$ с $c_s+c_c+c_h=1$
\begin{equation}
n=2c_s+c_c-c_h,\qquad K_2=6c_s+2c_c .
\end{equation}
Двух решений хватает, чтобы задать любое $n$, но тогда $K_2$ определяется через $n$. Например,
цилиндр и гиперболоид дают $n=2c_c-1$, $K_2=n+1$, и $n=0$ получается при $c_c=c_h=1/2$.
Три решения позволяют задать $n$ и $K_2$ независимо:
\begin{equation}
c_s=\frac{K_2-n-1}{3},\qquad c_c=n+1-\frac{K_2}{2},\qquad c_h=\frac{K_2}{6}-\frac{2n-1}{3}.
\end{equation}
Электроды --- эквипотенциали $\chi=\text{const}$ выбранной комбинации.

\section{Энергия и импульс в дефлекторе}

Сохранение полной энергии $\gamma mc^2+q\Phi=\gamma_{\rm out}mc^2$ даёт, как в главе,
\begin{equation}\label{eq:gloc}
\gamma=\gamma_{\rm out}(\delta)-\beta_s^2\gamma_s\chi ,
\end{equation}
и это равенство точно. Локальный импульс $P$ в единицах $p_s$:
\begin{equation}\label{eq:P}
\frac{P^2}{p_s^2}=\frac{\bigl(\gamma_{\rm out}(\delta)-\beta_s^2\gamma_s\chi\bigr)^2-1}{\beta_s^2\gamma_s^2}
=(1+\delta)^2-2\gamma_{\rm out}\frac{\chi}{\gamma_s}+\beta_s^2\chi^2 .
\end{equation}
В линейном порядке $P/p_s=1+\delta_{\rm loc}=1+\delta-\chi$, что совпадает с $\delta_{\rm loc}$ главы.

\section{Уравнение орбиты в дефлекторе}

Гамильтониан с $s$ в роли времени (поле статическое, $\delta$ сохраняется):
\begin{equation}\label{eq:H}
H=-\Bigl(1+\frac x\rho\Bigr)\sqrt{\frac{P^2(x,\delta)}{p_s^2}-p_x^2},
\end{equation}
$p_x$ --- поперечный импульс в единицах $p_s$. Из $x'=\partial H/\partial p_x$
\begin{equation}\label{eq:px}
p_x=x'\Bigl(1+\delta-2\frac x\rho\Bigr)+O(2),
\end{equation}
что даёт связь $p_x=x'-2xx'/\rho$ при $\delta=0$ на границе дефлектора. Исключая $p_x$ из
$x'=\partial H/\partial p_x$, $p_x'=-\partial H/\partial x$ и раскладывая до второго порядка
по $x$, $x'$, $\delta$, получаем
\begin{equation}\label{eq:orbit2}
\boxed{\begin{aligned}
x''+K\,x+A\,\frac{x^2}{\rho^3}
&=\Bigl(1+\frac1{\gamma_s^2}\Bigr)\frac{\delta}{\rho}+B\,\frac{x\delta}{\rho^2}+C\,\frac{\delta^2}{\rho}+\frac{x'^2}{\rho},
\end{aligned}}
\end{equation}
где
\begin{equation}\label{eq:coef}
\begin{gathered}
K=\frac{2-n+1/\gamma_s^2}{\rho^2},\qquad
A=\frac12\Bigl(2m-7n+8+\frac{10-3n}{\gamma_s^2}\Bigr),\\
B=4-n+\frac{8-n}{\gamma_s^2},\qquad
C=-\Bigl(1+\frac{5}{2\gamma_s^2}-\frac{1}{2\gamma_s^4}\Bigr).
\end{gathered}
\end{equation}
Линейная часть \eqref{eq:orbit2} --- уравнение $x''+x/\rho^2=1/\rho-1/R$ с $1/R=-qE_x/(pv)$ и
$pv/(p_sv_s)=1+(1+1/\gamma_s^2)(\delta-x/\rho)$, т.\,е. уравнение для $D_1$ главы.
Во втором порядке форма $x''+x/\rho^2=1/\rho-1/R$ уже неверна: нужно учесть удлинение
пути $ds=(1+x/\rho)ds_0$ и член $x'^2$, это и делает \eqref{eq:H}. В $A$ коэффициент $m$
входит через силу $qE_x$, $n$ --- через силу и через квадратичный член $\chi$ в $P^2$.

\section{Уравнения для $D_1$ и $D_2$}

Подставляем замкнутую орбиту $x=D_1\delta+D_2\delta^2$ в \eqref{eq:orbit2} и собираем члены
при $\delta$ и $\delta^2$:
\begin{equation}\label{eq:D1}
\boxed{D_1''+K\,D_1=\frac{1+1/\gamma_s^2}{\rho}}
\end{equation}
\begin{equation}\label{eq:D2}
\boxed{D_2''+K\,D_2=-A\,\frac{D_1^2}{\rho^3}+B\,\frac{D_1}{\rho^2}+\frac{C}{\rho}+\frac{D_1'^2}{\rho}}
\end{equation}
Уравнение для $D_2$ линейно с той же жёсткостью $K$, а правая часть известна после решения
\eqref{eq:D1}. Вне дефлекторов в магнитном квадруполе $x''+K_1x(1-\delta)=0$, и
\begin{equation}
D_1''+K_1D_1=0,\qquad D_2''+K_1D_2=K_1D_1 .
\end{equation}
На границах $D_2$ непрерывно, а $D_2'$ следует из непрерывности $p_x$: внутри
$p_x=x'(1+\delta-2x/\rho)$, снаружи $p_x=x'(1+\delta)$. Члены при $\delta^2$ дают на входе
\begin{equation}
D_2'^{\,\rm in}-D_2'^{\,\rm out}=\frac{2D_1D_1'}{\rho},
\end{equation}
на выходе --- с обратным знаком.

\paragraph{Проверка.} При $m=n(n+1)/2$, $n=1$, $\beta_s^2=1-1/\gamma_s^2$
\eqref{eq:orbit2} совпадает с уравнением второго порядка для цилиндрического дефлектора:
$A=(10-7\beta_s^2)/2$, $B=10-7\beta_s^2$, $C=-(6-3\beta_s^2-\beta_s^4)/2$.

\end{document}
